\subsubsection{Sublimation and condensation}
\label{subsec:sublimation_condensation}

Following Eq.~(36) of \citet{stammlerRedistributionCOLocation2017}, the exchange between
ice and vapor of species $X$ on grain population $k$ is written as
\begin{equation}
    \frac{\mathrm d\Sigma_{X^{\mathrm{ice}},k}}{\mathrm dt}
    =\chi_{X,k}\pi a_k^2 n_{\mathrm{dust},k}v_{\mathrm{therm},X}
    \left(\Sigma_{X^{\mathrm{vap}}}
    -\Sigma_{X^{\mathrm{vap}}}^{\mathrm{eq}}\right),
    \label{eq:simultaneous_phase_rate}
\end{equation}
where $\Sigma_{X^{\mathrm{vap}}}^{\mathrm{eq}}$ is the equilibrium
vapor surface density and $\chi_{X,k}$ indicates whether the grain
population participates in the exchange. All populations participate
in condensation, whereas sublimation requires an ice inventory
exceeding one molecular monolayer. The temperature, dust scale
heights, grain numbers, and material density are held fixed during
the chemistry step. Different radial cells exchange no material
through chemistry and can therefore be integrated independently.

The equilibrium vapor surface density of species $X$ is
\begin{equation}
    \Sigma_{X^{\mathrm{vap}}}^{\mathrm{eq}}
    =
    \sqrt{2\pi}H_g
    \frac{m_X}{k_{\mathrm B}T}P_X^{\mathrm{eq}}(T),
    \label{eq:sigma_vap_eq}
\end{equation}
where $P_X^{\mathrm{eq}}(T)$ is the equilibrium vapor pressure over ice
and $m_X$ is the molecular mass. The mean molecular thermal velocity is
\begin{equation}
    v_{\mathrm{therm},X}
    =\sqrt{\frac{8k_{\mathrm B}T}{\pi m_X}}.
\end{equation}

The midplane grain number density is
$n_{\mathrm{dust},k}=N_k/(\sqrt{2\pi}H_{\mathrm d,k})$
and remains fixed throughout chemistry.

For a fixed grain radius and a fixed set of participating populations,
we define
\begin{equation}
    C_{X,k}=\chi_{X,k}\pi a_k^2
    n_{\mathrm{dust},k}v_{\mathrm{therm},X},
    \qquad C_X=\sum_k C_{X,k}.
\end{equation}
Conservation of each volatile requires
\begin{equation}
    \frac{\mathrm d\Sigma_{X^{\mathrm{vap}}}}{\mathrm dt}
    =-\sum_k\frac{\mathrm d\Sigma_{X^{\mathrm{ice}},k}}{\mathrm dt}
    =-C_X\left(\Sigma_{X^{\mathrm{vap}}}
    -\Sigma_{X^{\mathrm{vap}}}^{\mathrm{eq}}\right).
\end{equation}
Thus, over an interval $h$ beginning at $t_n$, the vapor evolves as
\begin{equation}
    \Sigma_{X^{\mathrm{vap}}}(t_n+h)
    =\Sigma_{X^{\mathrm{vap}}}^{\mathrm{eq}}
    +\Delta_{X,n}\exp(-C_Xh),
    \label{eq:vapor_chem_solution}
\end{equation}
where $\Delta_{X,n}=\Sigma_{X^{\mathrm{vap}}}(t_n)
-\Sigma_{X^{\mathrm{vap}}}^{\mathrm{eq}}$. The corresponding ice
transfer is
\begin{equation}
    \delta\Sigma_{X^{\mathrm{ice}},k}
    =\frac{C_{X,k}}{C_X}\Delta_{X,n}
    \left[1-\exp(-C_Xh)\right],
    \label{eq:analytic_ice_substep}
\end{equation}
with zero transfer when $C_X=0$. Numerically, the exponential
difference is evaluated using \texttt{expm1} to avoid cancellation
for short intervals. Vapor is updated from the summed ice transfer,
so that phase exchange conserves the combined inventory of each
species to floating-point accuracy.

During sublimation, a population becomes inactive when its inventory
reaches
\begin{equation}
    \Sigma_{X^{\mathrm{ice}},k}^{\mathrm{monolayer}}
    =N_k\frac{4\pi(a_k+r_X)^2}{\sigma_X}m_X,
    \label{eq:monolayer_surface_density}
\end{equation}
where $N_k=\Sigma_{\mathrm d,k,0}/m_{k,0}$ is the fixed grain surface number density, $r_X$ and $m_X$
are the molecular radius and mass, and $\sigma_X$ is the area occupied
by one molecule. Writing $w_{X,k}=C_{X,k}/C_X$ and
$L_{X,k}=\Sigma_{X^{\mathrm{ice}},k}
-\Sigma_{X^{\mathrm{ice}},k}^{\mathrm{monolayer}}$, the depletion
time at fixed geometry is
\begin{equation}
    h_{X,k}^{\mathrm{dep}}
    =-\frac{1}{C_X}\ln\left(1-
    \frac{L_{X,k}}{w_{X,k}|\Delta_{X,n}|}\right).
\end{equation}
This expression applies when the fraction inside the logarithm lies
strictly between zero and one. Otherwise, depletion does not occur
in finite time before equilibrium is reached; populations already
at or below a monolayer are inactive. At the earliest depletion
event, the affected population is removed from the participating
set, the coefficients are recomputed, and the analytic solution is
continued over the remaining interval. Initially underfilled
populations are not raised to a monolayer.

Large phase changes can invalidate the fixed-radius approximation
over the full transport timestep. We therefore divide the chemistry
interval into adaptive local substeps. Let $\Sigma_{\mathrm d,k,0}$
and $a_{k,0}$ denote the total dust surface density and grain radius
at the start of chemistry. Since grain number and material density
remain fixed until remapping, the radius at a substep boundary is
\begin{equation}
    a_k(t_n)=a_{k,0}
    \left[1+\frac{\displaystyle\sum_X
    \left(\Sigma_{X^{\mathrm{ice}},k}(t_n)
    -\Sigma_{X^{\mathrm{ice}},k,0}\right)}
    {\Sigma_{\mathrm d,k,0}}\right]^{1/3}.
    \label{eq:simultaneous_radius_mass}
\end{equation}
All species use this same radius throughout a proposed substep.
Their transfers are calculated independently with
Eq.~\eqref{eq:analytic_ice_substep}, including depletion events,
and are then applied together. The shared radii and monolayer
inventories are recalculated before the next substep. This avoids
assigning different geometries to species according to their order
in the chemistry update.

To control the radius approximation, define the gross transferred
mass in population $k$ during a trial substep as
\begin{equation}
    Q_k=\sum_X\left|\delta\Sigma_{X^{\mathrm{ice}},k}\right|.
\end{equation}
We accept the substep only if every populated bin satisfies
\begin{equation}
    \frac{Q_k}{\Sigma_{\mathrm d,k}(t_n)}
    \leq 1-(1-\epsilon_a)^3,
    \label{eq:chemistry_radius_control}
\end{equation}
where $\Sigma_{\mathrm d,k}(t_n)$ includes all ice changes already
accepted during chemistry. Because each species transfers mass in
one direction at fixed temperature and geometry, $Q_k$ also bounds
the absolute net mass change at intermediate times in the trial
solution. Equation~\eqref{eq:chemistry_radius_control} therefore
bounds both growth and shrinkage of the reconstructed radius by
$\epsilon_a$. Using the gross transfer prevents condensation and
sublimation of different species from hiding a radius excursion by
canceling in the final net mass change.

Each cell first attempts its entire remaining chemistry interval.
If the trial violates Eq.~\eqref{eq:chemistry_radius_control}, it is
discarded and repeated with a shorter interval. Instantaneous phase
rates provide a timestep estimate, excluding species whose total
trial transfer is negligible compared with the permitted mass
change. This prevents a rapidly depleted, nearly bare species from
forcing unnecessarily short steps for other species. Acceptance
always depends on the actual trial transfers. Cells with little
phase exchange can complete chemistry in one substep, while cells
undergoing large mantle changes take more substeps without reducing
the global transport timestep. We adopt $\epsilon_a=0.01$ as the
default. This is a local geometry bound rather than a prescribed
relative error in the final ice inventories; convergence is assessed
by reducing $\epsilon_a$.

The final frozen-radius substep can leave an excess above the
monolayer corresponding to the final, smaller grains. We retain a
single additional shrinkage correction before remapping. It applies
to populations whose remaining inventory is at or below the
monolayer evaluated at the start of the full chemistry step. The
proposed release is the positive excess above the monolayer evaluated
at the final radius. It is accepted only if the sublimation rate,
evaluated using the vapor deficit after the proposed total release,
can remove it within $0.01\Delta t$. Corrections for all species are
calculated from the same state before being applied. Their additional
radius changes are not iterated, but their mass transfers are included
in the subsequent remapping. Remapping is performed once, after all
local chemistry substeps and this correction have completed.
